% Part: lambda-calculus
% Chapter: introduction
% Section: computable-lambda

\documentclass[../../../include/open-logic-section]{subfiles}

\begin{document}

\olfileid[id]{lam}{int}{lrp}
\olsection{Fungsi yang Dapat Dihitung Juga \usetoken{S}{lambda definable}}

\begin{thm}
\ollabel{thm:computable-lambda}
Setiap fungsi parsial yang dapat dihitung bersifat !!{lambda definable}.
\end{thm}

\begin{proof}
Kita perlu menunjukkan bahwa setiap fungsi parsial dapat dihitung~$f$
!!{lambda defined} oleh suatu suku lambda~$F$. Berdasarkan teorema bentuk
normal Kleene, cukup ditunjukkan bahwa setiap fungsi rekursif primitif
!!{lambda defined} oleh suatu suku lambda, kemudian bahwa fungsi-fungsi yang
!!{lambda definable} tertutup terhadap komposisi yang sesuai dan pencarian tak
terbatas. Untuk menunjukkan bahwa setiap fungsi rekursif primitif
!!{lambda defined} oleh suatu suku lambda, cukup ditunjukkan bahwa fungsi
awal bersifat !!{lambda definable}, serta bahwa fungsi parsial yang
!!{lambda definable} tertutup terhadap komposisi, rekursi primitif, dan
pencarian tak terbatas.
\end{proof}

Kita akan menggunakan notasi yang lebih konvensional agar sisa pembuktian lebih
mudah dibaca. Sebagai contoh, kita akan menulis $M(x,y,z)$ sebagai ganti
$Mxyz$. Meskipun notasi ini sugestif, ingatlah bahwa suku dalam kalkulus lambda
tanpa tipe tidak mempunyai aritas yang melekat; jadi, bagi suku~$M$ yang sama,
menulis $M(x,y)$ sama masuk akalnya dengan menulis $M(x,y,z,w)$. Namun,
penggunaan notasi ini menunjukkan bahwa kita memperlakukan $M$ sebagai fungsi
tiga variabel dan membantu memperjelas maksud definisinya. Dengan cara serupa,
kita akan mengatakan ``definisikan $M$ dengan $M(x,y,z)=\dots$'' sebagai ganti
``definisikan $M$ dengan $M=\lambd[x][\lambd[y][\lambd[z][\dots]]]$.''

\end{document}
